\documentclass[../../main.tex]{subfiles}
\begin{document}

\begin{flushright}
\footnotesize\textit{【自動文字起こし・要確認】}
\end{flushright}

{\bf [解]}

\paragraph{(1)}

$C_1$と$C_2$が共有点を持つ時，$y$を消去した
\begin{align}
-x^2+1&=(x-u)^2+u\\
2x^2-2ux+u^2+u-1&=0\quad\cdots①
\end{align}
が実解を持つ．判別式を$D$として，$D/4\ge0$だから，
\begin{align}
D/4=u^2-2(u^2+u-1)\ge0\iff-u^2-2u+2\ge0\iff-1-\sqrt{3}\le u\le-1+\sqrt{3}
\end{align}
となり，$(a,b)=(-1-\sqrt{3},-1+\sqrt{3})$．

\paragraph{(2)}

①の2解を$x_1,x_2$とする．対称性から$x_1\le x_2$としてよい．
\begin{align}
y_k=-x_k^2+1\quad(k=1,2)
\end{align}
より，与式を$f(u)$として，
\begin{align}
\dfrac{1}{2}f(u)&=\left|x_1(1-x_2^2)-x_2(1-x_1^2)\right|\\
&=\left|x_1-x_1x_2^2-x_2+x_1^2x_2\right|\\
&=\left|(x_1-x_2)+(x_1-x_2)x_1x_2\right|\\
&=|x_2-x_1||1+x_1x_2|\quad\cdots②
\end{align}

ここで，$x_1\le x_2$だから，①の解と係数の関係より
\begin{align}
x_2-x_1=\sqrt{-u^2-2u+2}
\end{align}
であり，また$x_1x_2=\dfrac{1}{2}(u^2+u-1)$だから，以上を②に代入して，
\begin{align}
f(u)&=2\sqrt{-u^2-2u+2}\left(1+\dfrac{u^2+u-1}{2}\right)\\
&=\sqrt{-u^2-2u+2}\,(u^2+u+1)
\end{align}

（$u^2+u+1$の判別式は$1-4<0$だから常に正で，絶対値記号を外せる．）

\paragraph{(3)}

$t=u+1$とおく．
\begin{align}
f(u)&=\sqrt{-(t-1)^2+3}\,\left((t-1)^2+(t-1)+1\right)\\
&=\sqrt{-t^2+3}\,(t^2-t+1)
\end{align}
であるから，(1)より，
\begin{align}
I=\int_{-\sqrt{3}}^{\sqrt{3}}\sqrt{-t^2+3}\,(t^2-t+1)\,dt=2\int_0^{\sqrt{3}}\sqrt{-t^2+3}\,(t^2+1)\,dt\quad\cdots③
\end{align}
である（奇関数の項$-t\sqrt{-t^2+3}$の積分は$0$になることを用いた）．各項を計算する．まず第一項は，$t=\sqrt{3}\sin\theta$とおく．$\dfrac{dt}{d\theta}=\sqrt{3}\cos\theta$だから，
\begin{align}
\int_0^{\sqrt{3}}t^2\sqrt{-t^2+3}\,dt&=\int_0^{\pi/2}3\sin^2\theta\cdot\sqrt{3}\cos\theta\cdot\sqrt{3}\cos\theta\,d\theta\\
&=9\int_0^{\pi/2}\dfrac{1}{4}\sin^22\theta\,d\theta\\
&=\dfrac{9}{4}\cdot\dfrac{1}{2}\left[\theta-\dfrac{1}{4}\sin4\theta\right]_0^{\pi/2}=\dfrac{9}{16}\pi
\end{align}

次項は，半径$\sqrt{3}$の四分円の面積に等しく，
\begin{align}
\int_0^{\sqrt{3}}\sqrt{-t^2+3}\,dt=\dfrac{1}{4}\pi(\sqrt{3})^2=\dfrac{3}{4}\pi
\end{align}

以上を③に代入して
\begin{align}
I=2\left(\dfrac{9}{16}\pi+\dfrac{3}{4}\pi\right)=\dfrac{21}{8}\pi
\end{align}

\end{document}
